\input{../preamble.tex}

\SetDate[14/04/2026]
\reversemarginpar
\setlength{\marginparsep}{.5cm}

\begin{document}
\pagestyle{fancy}
\fancyhead[L]{Seconde 5}
\fancyhead[C]{\textbf{Évaluation — Probabilités et Fonctions affines}}
\fancyhead[R]{\today}

\null\vspace{-20pt}
Consignes particulières : 
\begin{itemize}[label=$\bullet$]
	\item 
	La calculatrice est {autorisée}.
	\item
	Sauf mention du contraire, toutes les réponses doivent être justifiées.
	\item
	L'exercice \ref{exe:fa1} peut être entièrement fait sur la feuille d'évaluation.
	\item
	L'évaluation fait \pageref{lastpage} pages. La somme des points est \total{points}.
\end{itemize}

\marginpar{[pts]}
\hrule


\section*{Probabilités}


\exe{4}{
	On considère le lancer d'un D10, dé à 10 faces numérotées de $1$ à $10$.
	On note le nombre de la face du dessus.
	
	Soient $A, B$ les deux événements suivants.
		\begin{align*}
			A : \text{\og le nombre est pair \fg}, && B : \text{\og le nombre est multiple de $5$ \fg}.
		\end{align*}
		
	\begin{enumerate}
		\item Donner l'univers de l'expérience. % ainsi que son cardinal $|\Omega|$.
		\item Donner le cardinal de l'univers.
		%\item Donner l'ensemble des issues associé à l'événement $B$.
		\item Décrire avec des mots l'événement $A \cap B$. % puis donner l'ensemble des issues qui lui est associé.
		\item Décrire avec des mots l'événement $A \cup B$. %puis donner l'ensemble des issues qui lui est associé.
		%\item Décrire avec des mots l'événement $\overline{A}$. %puis donner l'ensemble des issues qui lui est associé.
	\end{enumerate}
}{exe:p1}{
	\begin{enumerate}
		\item 
		$\Omega = \bigset{ 1 ; 2 ; 3 ; 4 ; 5 ; 6 ; 7 ;8 ; 9 ; 10}$.
		\item 
		$|\Omega| = 10$.
		\item 
		$A\cap B$ : « le nombre est pair \textbf{et} multiple de 5 ».
		\item 
		$A\cup B$ : « le nombre est pair \textbf{ou} multiple de 5 ».
	\end{enumerate}
}


\exe{6, difficulty=1}{
	On tire une boule dans une urne contenant $3$ boules rouges et $4$ boules vertes et on regarde sa couleur.
	\begin{enumerate}[label=$\bullet$]
		\item Si on tire une boule rouge, on lance un D$3$ bien équilibré, dé à $3$ faces numérotées de $1$ à $3$, et on note la face du dessus.
		\item Si on tire une boule verte, on lance un D$5$ bien équilibré, dé à $5$ faces numérotées de $1$ à $5$, et on note la face du dessus.
	\end{enumerate}
	À la fin de l'expérience, on remet la boule tirée dans l'urne.

	\begin{enumerate}
		\item 
		Créer un arbre de probabilité pour cette situation.
		\item 
		Quelle est la probabilité \emph{exacte} d'obtenir 4 ? 
		\item
		Considérons l'événement $E$ : « obtenir au moins un 4 après dix tirages consécutifs ».
			\begin{enumerate}
				\item Décrire $\overline{E}$ avec des mots.
				\item Calculer $P(E)$ à $10^{-2}$ près.
			\end{enumerate}
	\end{enumerate}

}{exe:p2}{
	\begin{enumerate}
		\item \,\\
		\begin{center}
		\begin{tikzpicture}
			% depth 1
			\foreach \i in {-3, 3}
			\draw[-, thick, black] (0,0) node {$\bullet$} -- (\i,-2);
			% depth 2
			% i = -3
			\foreach \foreach \j in {-1, 0, 1}
				\draw[-, thick, black] (-3,-2) node {$\bullet$} -- (-3+\j,-4) node {$\bullet$};
			% i = +3
			\foreach \j in {-2, -1, 0, 1, 2}
				\draw[-, thick, black] (3,-2) node {$\bullet$} -- (3+\j,-4) node {$\bullet$};
			
			\draw (-3,-2) node[above left] {Rouge};
			\draw (3,-2) node[above right] {Vert};
				
			\draw (-4,-4) node[below] {1};
			\draw (-3,-4) node[below] {2};
			\draw (-2,-4) node[below] {3};
			
			\draw (1,-4) node[below] {1};
			\draw (2,-4) node[below] {2};
			\draw (3,-4) node[below] {3};
			\draw (4,-4) node[below] {4};
			\draw (5,-4) node[below] {5};
			
			% proba
			\draw (-1.75,-.5) node {$\dfrac37$};
			\draw (1.75,-.5) node {$\dfrac47$};
			
			\draw (-2,-2.5) node {$\frac13$};
			\draw (-2.75,-3.25) node {$\frac13$};
			\draw (-4,-2.5) node {$\frac13$};
			
			\draw (2,-2.5) node {$\frac15$};
			\draw (2.25,-3.1) node {$\frac15$};
			\draw (3.2,-3.25) node {$\frac15$};
			\draw (3.75,-3.1) node {$\frac15$};
			\draw (4.25,-2.5) node {$\frac15$};
			
			% red path
			\draw[dashed, thick, PINK, xshift=3pt] (0,0) -- (3,-2) -- (4,-4);
		\end{tikzpicture}
		\end{center}
		\item 
		Seul le chemin en pointillés roses aboutit à obtenir 4 : sa probabilité est
			\[ \frac47 \times \frac15 = \frac4{35}. \]
		\item
			\begin{enumerate}
				\item 
				$\overline{E}$ : « n'obtenir aucun 4 après dix tirages consécutifs ».
				\item 
				D'abord, 
					\begin{align*}
					P\bigl(\overline{E}\bigr) &= P(\text{« ne pas obtenir 4 »})^{10} \\
						&= \left(1 - \frac4{35}\right)^{10} \\
						&= \left(\frac{31}{35}\right)^{10}
					\end{align*}
				Par conséquent, 
					\[ P(E) = 1 - P\bigl(\overline{E}\bigr) = 1 - \left(\frac{31}{35}\right)^{10} \approx 0,70. \]
			\end{enumerate}
	\end{enumerate}
}

\vspace{10pt}

\hrule

\vspace{10pt}
%Brouillon pour Natalija

\newpage

\section*{Fonctions affines}


\begin{thm}\label{thm:1}
	Soit $f(x) = ax + b$ une fonction affine sur $\R$.
	Considérons $A(x_A ; y_A)$ et $B(x_B ; y_B)$ deux points distincts appartenant à $\C_f$.
	
	Alors le coefficient directeur $a$ est égal à
		\[ a = \phantom{\frac{y_B-y_A}{x_B - x_A} \hspace{1cm}} \]
\end{thm}

\exe{2}{
	Compléter le théorème \ref{thm:1} vu en classe.
}{exe:fa1}{
	\[ a = \dfrac{y_B - y_A}{x_B - y_A} = \dfrac{y_A - y_B}{x_A - y_B}. \]
}

%\exe{2}{
%	Soit $F$ la fonction affine sur $\R$ donnée par
%		\[ F(x) = 6+x. \]
%	Donner sans justification le coefficient directeur $a$ et l'ordonnée à l'origine $b$ de $F$.
%}{exe:fa4}{
%
%}

\exe{2}{
	Soit $f$ la fonction affine sur $\R$ donnée par
		\[ f(x) = x+31. \]
	Déterminer l'expression algébrique de la fonction affine $g$ telle que $\C_g$ soit parallèle à $\C_f$ et que $\C_g$ passe par $(3;-1)$.
}{exe:fa2}{
	Posons $g(x) = ax+b$ avec $a$ et $b$ deux paramètres à trouver à l'aide des informations disponibles.
	
	Premièrement, le parallélisme implique que $a$ est égal au coefficient directeur de $f$, qui est 1.
	Donc $a=1$ et $g(x) = x+b$.
	
	Deuxièmement, l'appartenance de $(3;-1)$ à $\C_g$ implique que $g(3) = -1$ et donc que $3+b = -1$, équation de laquelle nous obtenons $b=-4$.
	
	En conclusion, $g(x) = x-4$.
}

\exe{3, difficulty=1}{
	Pour chaque paire de fonctions affines $f, g$ sur $\R$, donner $\C_f \cap \C_g$.
	\begin{multicols}{2}
	\begin{enumerate}
		\item $f(x) = \dfrac53 x - 4$ et $g(x) = 11$.
		\item $f(x) = 2x+1$ et $g(x) = 2x-3$.
	\end{enumerate}
	\end{multicols}
}{exe:fa3}{

	\begin{enumerate}
		\item 
		En posant $f(x) = g(x)$, nous obtenons l'abscisse du point d'intersection des droites $\C_f$ et $\C_g$.
			\begin{align*}
				f(x) &= g(x) \\
				\dfrac53 x - 4 &= 11 \\
				\dfrac53 x &= 15 \\
				x &= 15 \cdot \frac35 \\
				x &= 9
			\end{align*}
		L'ordonnée du point d'intersection est égale à $y = f(9) = g(9) = 11$, donc
			\[ \C_f \cap \C_g = \bigset{ (9 ; 11) }. \]
		\item 
		Les deux fonctions admettent le même coefficient directeur mais deux ordonnées à l'origine différentes : les droites sont donc parallèles non confondues et
			\[ \C_f \cap \C_g = \emptyset. \]
	\end{enumerate}
}

\exe{2}{
	Considérons une fonction affine $f$ telle que les points $(2;6)$ et $(-1 ; 2)$ appartiennent à $\C_f$.
	
	\underline{Sans calculer} l'expression algébrique de $f$, répondre aux questions suivantes.
	\begin{enumerate}
		\item Tracer $\C_f$ sur le domaine $[-3; 3]$.
		\item Estimer \emph{graphiquement} à l'aide de $\C_f$ la valeur de l'ordonnée à l'origine de $f$.
	\end{enumerate}
}{exe:fa5}{
	Il faut utiliser ici que $\C_f$ est une droite passant par les deux points donnés.
	En plaçant les points et en traçant la droite, l'ordonnée à l'origine $b$ se lit comme l'ordonnée du point d'abscisse nulle, soit $(0 ; b)$.
	
	\begin{center}
	\begin{tikzpicture}
	\begin{axis}[xmin = -3, xmax=3, ymin=-1, ymax=8, axis x line=middle, axis y line=middle, axis line style=<->, xlabel={}, ylabel={}, grid=both, ytick distance = 1]
		\addplot[mark=*, mark size=1pt, PINK] (2,6) node[above]{$(2;6)$};
		\addplot[mark=*, mark size=1pt, ORANGE] (-1,2) node[below]{$(-1;2)$};
		
		\addplot[-, thick, BLUE_E] (2,6) -- (axis cs:-1,2) node[midway,above]{$\C_f$};
		
		\addplot[mark=*, mark size=1pt, GREEN_E] (0,3.333) node[below right]{$(0;3,3)$};
	\end{axis}
	\end{tikzpicture}
	\end{center}
	
	Il suit que $b\approx 3,3$.
}

\exe{2, difficulty=1}{
	\mbox{Considérons une fonction affine $f$ telle que les points $A(2;6)$ et $B(-1 ; 2)$ appartiennent à $\C_f$.}
	
	Trouver $f$ \emph{exactement} en calculant son coefficient directeur puis son ordonnée à l'origine.
}{exe:fa6}{
	Posons $f(x) = ax+b$ avec $a$ et $b$ des paramètres à déterminer à l'aide de l'appartenance des points $A$ et $B$ à la courbe $\C_f$.
	
	Premièrement, d'après le théorème \ref{thm:1}, nous avons
		\[ a = \dfrac{y_B - y_A}{x_B-x_A} = \frac{2-6}{-1-2} = \frac{-4}{-3} = \frac43. \]
	D'où $f(x) = \frac43 x + b$.
	Ceci est cohérent avec l'exercice \ref{exe:fa5} car lorsque $x$ augmente de 3 pour passer de $B$ à $A$, $y$ augmente de 4.
	Donc par proportionnalité, quand $x$ augmente de 1, $y$ augmente de $\frac43$, le coefficient directeur.
	
	Deuxièmement, l'appartenance de $A$ à $\C_f$ donne
		\begin{align*}
			f(x_A) &= y_A \\
			f(2) &= 6 \\
			\frac43\times2 + b &= 6 \\
			b &= 6 - \frac83 = \frac{10}3
		\end{align*}
	En conclusion, $f(x) = \frac43x + \frac{10}3$.
	Ceci est cohérent avec le $b\approx3,3$ trouvé à l'exercice \ref{exe:fa5}, car $\frac{10}3 = 3 + \frac13 = 3,333\hdots$.
}

\vspace{10pt}

\hrule

\vspace{10pt}
%Brouillon pour Natalija

%%%%%%%%%%%

\label{lastpage}
\newpage
\fancyhead[C]{\textbf{Solutions}}
\shipoutAnswer
	
\end{document}
